% ============================================================
%  01_introduction.tex — Section 1 : Introduction et objectif
% ============================================================
\section{Introduction et objectif}

\subsection{Contexte}

Ce document décrit l'implémentation concrète du solveur CSP (Constraint Satisfaction Problem)
pour la génération de structures polycycliques non-benzenoïdes. Il complète le document de
modélisation formelle (\textit{Stage\_M2\_IMD\_modelisation\_csp.tex}) qui définit les contraintes
mathématiques, et le document de validation (\textit{validation\_reconstruction\_csp.tex}) qui
traite de la reconstruction 3D.

\subsection{Le problème}

Étant donné un \textbf{benzénoïde} (assemblage plan d'hexagones fusionnés par arête), on cherche
toutes les façons de remplacer certains hexagones par des \textbf{pentagones} (5 côtés) ou des
\textbf{heptagones} (7 côtés), sous les contraintes suivantes :
\begin{enumerate}
  \item Le nombre total de carbones est conservé (contrainte d'Euler).
  \item Chaque voisinage local est chimiquement réalisable (table de voisinage).
  \item Deux solutions structurellement identiques ne sont comptées qu'une fois (symétrie).
\end{enumerate}

La \textbf{table de voisinage} est issue du générateur (cf. \textit{documentation.tex}) :
1209 configurations locales validées par optimisation xTB (GFN2-xTB) et test de planarité
(seuil 10°), dont 6 configurations purement benzenoïdes.
